\begin{figure}%[h]
    \centering
    \begin{tabular}{
        M{0.0\linewidth}@{\hspace{0\tabcolsep}}
        M{0.45\linewidth}@{\hspace{0\tabcolsep}}
        M{0.45\linewidth}@{\hspace{0\tabcolsep}}
        M{0.05\linewidth}@{\hspace{0\tabcolsep}}
    }
        \rotatebox[x=7pt,y=-12pt]{90}{\tiny $x_2$} &
        \begin{tabular}{
            M{0.25\linewidth}@{\hspace{0\tabcolsep}}
            M{0.25\linewidth}@{\hspace{0\tabcolsep}}
            M{0.25\linewidth}@{\hspace{0\tabcolsep}}
            M{0.25\linewidth}@{\hspace{0\tabcolsep}}
        }
            \algo & \ddrm & \dps & \rnvp \\
            \includegraphics[width=\hsize]{images/optimal/gaussian/800_1_1_25_20_MCG_DIFF}
            & \includegraphics[width=\hsize]{images/optimal/gaussian/800_1_1_25_20_DDRM}
            & \includegraphics[width=\hsize]{images/optimal/gaussian/800_1_1_25_20_DPS}
            & \includegraphics[width=\hsize]{images/optimal/gaussian/800_1_1_25_20_RNVP} \\
        \end{tabular}
    &
        \begin{tabular}{
                M{0.25\linewidth}@{\hspace{0\tabcolsep}}
                M{0.25\linewidth}@{\hspace{0\tabcolsep}}
                M{0.25\linewidth}@{\hspace{0\tabcolsep}}
                M{0.25\linewidth}@{\hspace{0\tabcolsep}}
            }
                  \algo & \ddrm & \dps & \rnvp \\
                  \includegraphics[width=\hsize]{images/optimal/funnel/10_1_1_20_19_MCG_DIFF}
                  & \includegraphics[width=\hsize]{images/optimal/funnel/10_1_1_20_19_DDRM}
                  & \includegraphics[width=\hsize]{images/optimal/funnel/10_1_1_20_19_DPS}
                  & \includegraphics[width=\hsize]{images/optimal/funnel/10_1_1_20_19_RNVP}
        \end{tabular}
        & \rotatebox[x=10pt,y=-15pt]{-90}{\tiny $\operatorname{PCA}_2$}\\\addlinespace[-1pt]
    & \tiny $x_1$ & \tiny $\operatorname{PCA}_1$ &
    \end{tabular}
    \caption{The first and last four columns correspond respectively to GM with $(\dimx, \dimtr) = (800, 1)$ and FM with $(\dimx, \dimtr) = (10, 1)$.
    The blue and red dots represent respectively samples from the exact posterior and those generated by each of the algorithms used (names on top).}
    \label{fig:gmm:20:illustration}
\end{figure}
\DTLloaddb{comparison_gmm_sw_20}{data/gaussian/main_25_20_inverse_problem_comparison.csv}
\DTLloaddb{comparison_fmm_sw_100}{data/funnel/main_20_100_inverse_problem_comparison.csv}
\begin{table}%[h]
    \centering
    \begin{subtable}[h]{0.45\textwidth}
        \resizebox{1\textwidth}{!}{
            \begin{tabular}{|c|c|c|c|c|c|c|}%
                \hline
                $d$ & $d_y$ & \algo & \ddrm & \dps & \rnvp \\
                \hline
                \DTLforeach*{comparison_gmm_sw_20}{
                    \dim=D_X,
                    \ydim=D_Y,
                    \distours=MCG_DIFF,
                    \distoursnum=MCG_DIFF_num,
                    \distddrm=DDRM,
                    \distddrmnum=DDRM_num,
                    \distdps=DPS,
                    \distdpsnum=DPS_num,
                    \distrnvp=RNVP,
                    \distrnvpnum=RNVP_num}{
                    \dim & \ydim &
                    \DTLgminall{\rowmin}{\distoursnum,\distddrmnum,\distdpsnum,\distrnvpnum}%
                    \dtlifnumeq{\distoursnum}{\rowmin}{\bf}{}\distours &
                    \dtlifnumeq{\distddrmnum}{\rowmin}{\bf}{}\distddrm &
                    \dtlifnumeq{\distdpsnum}{\rowmin}{\bf}{}\distdps &
                    \dtlifnumeq{\distrnvpnum}{\rowmin}{\bf}{}\distrnvp\DTLiflastrow{}{\\
                    }}
                \\\hline
            \end{tabular}%
        }
    \end{subtable}%
    \begin{subtable}[h]{0.45\textwidth}
        \resizebox{1\textwidth}{!}{
            \begin{tabular}{|c|c|c|c|c|c|c|}%
                \hline
                $d$ & $d_y$ & \algo & \ddrm & \dps & \rnvp \\
                \hline
                \DTLforeach*{comparison_fmm_sw_100}{
                    \dim=D_X,
                    \ydim=D_Y,
                    \distours=MCG_DIFF,
                    \distoursnum=MCG_DIFF_num,
                    \distddrm=DDRM,
                    \distddrmnum=DDRM_num,
                    \distdps=DPS,
                    \distdpsnum=DPS_num,
                    \distrnvp=RNVP,
                    \distrnvpnum=RNVP_num}{
                    \dim & \ydim &
                    \DTLgminall{\rowmin}{\distoursnum,\distddrmnum,\distdpsnum,\distrnvpnum}%
                    \dtlifnumeq{\distoursnum}{\rowmin}{\bf}{}\distours &
                    \dtlifnumeq{\distddrmnum}{\rowmin}{\bf}{}\distddrm &
                    \dtlifnumeq{\distdpsnum}{\rowmin}{\bf}{}\distdps &
                    \dtlifnumeq{\distrnvpnum}{\rowmin}{\bf}{}\distrnvp\DTLiflastrow{}{\\
                    }}
                \\\hline
            \end{tabular}%
        }
    \end{subtable}%
    \caption{Sliced Wasserstein for the GM (left) and FM (right) case.}
    \label{table:gmm:sw}
\end{table}
\vspace{-.3mm}
We refer to the Funnel mixture prior as FM prior (see \cref{appendix:numerics} for the definition).
For GM prior, we consider a mixture of $25$ components with pre-established means and variances.
For FM prior, we consider a mixture of $20$ components consisting of rotated and translated funnel distributions.
For a given pair $(\dimx, \dimtr)$, we sample a prior distribution by randomly sampling the weights of the mixture and
for the FM case the translation and rotation of each component.
We then randomly sample measurement models $(\lmeas, \m{A}, \noisestd) \in \rset^{\dimtr} \times \mset{\dimtr}{\dimx} \times [0, 1]$.
For each pair of prior distribution and measurement model, we generate $10^4$ samples from \algo, \dps, \ddrm, \rnvp,
and from the posterior either analytically (GM) or using NUTS (FM).
We calculate for each algorithm the sliced Wasserstein (SW) distance between the resulting samples and the posterior samples.
\Cref{table:gmm:sw} shows the CLT $95\%$ confidence intervals obtained over $20$ seeds.
\Cref{fig:gmm:20:illustration} illustrate the samples for the different algorithms for a given seed.
We see that \algo\ outperforms all the other algorithms in each setting tested.
The complete details of the numerical experiments performed in this section is available in \cref{appendix:numerics} as well as an additional visualisations.